\documentclass{article}

\usepackage{amsmath}

\title{Beverton--Holt Population Model}
\author{}
\date{\today}

\begin{document}

\maketitle
\section{Beverton-Holt model}
Let $N_t \geq 0$ denote the population at time $t$. We use the
scaled Beverton--Holt model
$$
    N_{t+1} = F(N_t)
            = \frac{K r N_t}{K + N_t},
$$
where $r \geq 0$ is the growth parameter and $K>0$ is the
population-scale parameter.

\subsection{Equilibria}

An equilibrium $N^*$ satisfies $F(N^*)=N^*$. Therefore,
$$
    N^* = \frac{K r N^*}{K + N^*}.
$$
One solution is the extinction equilibrium
$$
    N^*_0=0.
$$
For a positive equilibrium, $N^*>0$, we divide by $N^*$ and obtain
\begin{align*}
    1 &= \frac{K r}{K+N^*},\\
    K+N^* &= Kr,\\
    N^* &= K(r-1).
\end{align*}

The positive equilibrium exists
only when $r>1$ and the population can not be negative. Thus,
$$
\begin{cases}
    r\leq 1: & N^*=0 \text{ is the only equilibrium},\\
    r>1: & N^*=0 \text{ and } N^*=K(r-1) \text{ are equilibria}.
\end{cases}
$$

\subsection{Stability}

The derivative of the population map is
$$
    F'(N)=\frac{K^2r}{(K+N)^2}.
$$
At extinction,
$$
    F'(0)=r.
$$
Therefore, extinction is locally stable when $0\leq r<1$, unstable
when $r>1$, and is a threshold case when $r=1$.

At the positive equilibrium,
$$
    F'\bigl(K(r-1)\bigr)=\frac{1}{r}.
$$
Consequently, whenever the positive equilibrium exists, $r>1$, it is
locally stable.

\section{Introduction of Fishing Effort}

We now introduce a fishing effort parameter $e$, where
$$
    0 \leq e \leq 1.
$$
The value $e=0$ means no fishing, while $e=1$ means complete removal
after reproduction. With fishing, the one-population model becomes
$$
    N_{t+1}
    = F_e(N_t)
    = (1-e)\frac{K r N_t}{K+N_t}.
$$

An equilibrium satisfies $F_e(N^*)=N^*$:
$$
    N^* = (1-e)\frac{K r N^*}{K+N^*}.
$$
Again, $N^*=0$ is always an equilibrium. For a positive equilibrium,
$N^*>0$, we divide by $N^*$ and obtain
\begin{align*}
    1 &= (1-e)\frac{K r}{K+N^*},\\
    K+N^* &= K r(1-e),\\
    N^* &= K\bigl(r(1-e)-1\bigr).
\end{align*}

Therefore, the positive equilibrium exists only when
$$
    r(1-e)>1.
$$
Equivalently, if $r>1$, persistence requires
$$
    e < 1-\frac{1}{r}.
$$
If $r(1-e)\leq 1$, then extinction is the only nonnegative equilibrium.

The derivative of the fishing map is
$$
    F_e'(N)=\frac{K^2r(1-e)}{(K+N)^2}.
$$
At extinction,
$$
    F_e'(0)=r(1-e).
$$
Thus, extinction is locally stable when $r(1-e)<1$ and unstable when
$r(1-e)>1$.

At the positive equilibrium,
$$
    F_e'\bigl(K(r(1-e)-1)\bigr)=\frac{1}{r(1-e)}.
$$
Whenever the positive equilibrium exists, $r(1-e)>1$, this derivative
has absolute value less than $1$, so the positive equilibrium is locally
stable.

\end{document}
